\chapter{Anexo}
% \label{anexo:revisao}

% \begin{table}[!htb]
% \caption{Extração de resultados pela Google Scholar}
% \begin{center}
% {
% \footnotesize
% \resizebox{\textwidth}{!}{
% \begin{tabular}{
% |p{\linewidth}|
% }
% \hline
% \cellcolor{gray!15}
% \textbf{Pontos relevantes para nosso trabalho} \\
% \hline
% Arquitetura de microsserviços, Kubernetes para orquestração, Microsserviços em contêineres, Tratamento de disponibilidade de serviços, Medição de disponibilidade de serviços como porcentagem do tempo que o serviço é provisionado, Avaliação do Kubernetes como provedor de disponibilidade de serviços, Mostraram que o Kubernetes não satisfaz a disponibilidade de serviço de 99,999\% do tempo, e em alguns casos falha em garantir recuperação de sistema, propõem "HA State Controller" integrado com o Kubernetes permitindo replicação de estados e redirecionamento de serviços automático. Os resultados das investigações mostraram ganhos em tempo de recuperação do sistema em 50\% \cite{VAYGHAN2021110924}. \\
% \hline





% \\
% \hline\\
% \hline\\
% \hline\\
% \hline\\
% \hline\\
% \hline\\
% \hline\\
% \hline


% \end{tabular}
% }
% }
% \end{center}
% \label{tab:extracao-GoogleScholar}
% \end{table}
